% ID: SEL-D0059-OB-E03
% ESTADO: cotejada. Segunda pasada visual y borrador compilado revisados.
% Fuente: Selectividad/sel_2011_matematicas_aplicadas/222-EXAMEN 2.pdf
% SHA-256: 998f20402054ca6d0454da83ffc8e1ff6951a576b48836a32e81e524c5985214
% Página: 2; opción B; ejercicio original 3.

Una bolsa contiene 5 bolas blancas, 3 rojas y 4 negras. Ana y Manolo practican el siguiente juego: Ana saca una bola, anota su color y la devuelve a la bolsa, a continuación Manolo extrae una bola y anota su color. Si las dos bolas extraídas tienen el mismo color gana Ana, si sólo hay una bola blanca gana Manolo, y en otro caso hay empate.
\begin{enumerate}[label=\alph*)]
\item \textbf{(1.25 puntos)} Calcule la probabilidad de que gane Ana.
\item \textbf{(1 punto)} Calcule la probabilidad de que gane Manolo.
\item \textbf{(0.25 puntos)} Calcule la probabilidad de que haya empate.
\end{enumerate}
